\documentclass[11pt,a4paper]{article}
\usepackage[utf8]{inputenc}
\usepackage[T1]{fontenc}
\usepackage{amsmath,amssymb}
\usepackage{graphicx}
\usepackage{booktabs}
\usepackage{hyperref}
\usepackage[margin=2.5cm]{geometry}
\usepackage{natbib}
\usepackage{float}

\title{Gaia: A Soil-Specialized Foundation Model for Microbiome Understanding and Agricultural Prediction}

\author{
  Hajin Nam \\
  Department of Mathematics, College of Natural Sciences \\
  Keimyung University, Daegu, South Korea \\
  \texttt{hajin0127@gmail.com} \\
  \texttt{https://github.com/Kimchikilla/Gaia}
}

\date{}

\begin{document}
\maketitle

\begin{abstract}
We present Gaia, a foundation model specialized for soil microbiome analysis, built through continual pre-training of MGM (Microbiome General Model) on 8,329 soil samples from MGnify, Earth Microbiome Project, and experimental datasets. Gaia achieves near-parity with Random Forest on soil biome classification (98.6\% vs 98.7\%) and drought detection (91.9\% vs 92.0\%), while demonstrating superior performance on temporal yield prediction ($R^2$=0.541 vs 0.504). Additionally, Gaia predicts soil total carbon ($R^2$=0.829) and pH ($R^2$=0.947) from microbiome composition alone. Functional profile analysis reveals that drought conditions increase carbon decomposition activity by 98.6\% while decreasing nitrogen fixation by 13.4\%. We expand the microbial vocabulary from 9,669 to 12,916 genera through a two-step training strategy that preserves pre-trained knowledge while incorporating soil-specific taxa. Across ten benchmark tasks, Gaia demonstrates that soil-specialized foundation models can match or exceed traditional machine learning approaches, with particular advantages in temporal prediction and functional analysis. All code, data pipelines, and model weights are publicly available at \url{https://github.com/Kimchikilla/Gaia}.
\end{abstract}

\section{Introduction}

Soil microbiomes play a critical role in ecosystem functioning, agricultural productivity, and global biogeochemical cycles. The soil harbors an estimated $10^9$ microbial cells per gram, forming complex communities whose composition reflects and influences soil health, nutrient cycling, and plant growth. Despite advances in high-throughput sequencing, interpreting these complex communities remains challenging.

Traditional approaches to soil microbiome analysis rely on statistical methods such as Random Forest, which treat each microbial taxon independently. These methods cannot capture the intricate ecological relationships between microorganisms---co-occurrence patterns, metabolic dependencies, and competitive interactions that collectively determine community function.

Foundation models, pre-trained on large corpora through self-supervised learning, have revolutionized natural language processing and protein structure prediction. Recently, MGM (Microbiome General Model) \citep{mgm2024} demonstrated that transformer-based language models can learn meaningful representations of microbial communities by treating abundance-ranked genus sequences as ``sentences'' in a microbial ``language.'' Concurrently, BiomeGPT applied a similar approach to human gut microbiome disease classification.

However, these models were trained on all biomes or specific human-associated environments without soil specialization. Soil microbiomes have distinct characteristics: exceptionally high diversity (typically 100--500+ genera per sample), strong environmental filtering by pH and moisture, seasonal dynamics tied to agricultural practices, and direct relevance to food security and climate change mitigation through carbon sequestration.

We introduce Gaia, a soil-specialized foundation model that addresses these limitations through:
\begin{enumerate}
\item Continual pre-training on 8,329 soil-specific samples from four public databases
\item Vocabulary expansion from 9,669 to 12,916 genera to cover soil-specific taxa
\item A two-step training strategy that preserves pre-trained knowledge
\item Comprehensive evaluation across ten benchmark tasks including temporal prediction and carbon estimation
\item Functional profile analysis linking microbial communities to ecological functions
\end{enumerate}

\section{Methods}

\subsection{Data Collection}

We assembled soil microbiome data from four public sources (Table~\ref{tab:data}):

\begin{table}[H]
\centering
\caption{Data sources for Gaia training}
\label{tab:data}
\begin{tabular}{lrrl}
\toprule
Source & Samples & Genera & Description \\
\midrule
MGnify & 2,790 & 2,643 & API-collected soil biome samples \\
Earth Microbiome Project & 4,628 & 1,659 & Global standardized soil samples \\
Naylor et al. & 623 & 1,001 & Drought stress experiment (16S, DADA2) \\
BonaRes Westerfeld & 288 & 1,864 & 20-year tillage/fertilization trial \\
\midrule
\textbf{Total} & \textbf{8,329} & & \\
\bottomrule
\end{tabular}
\end{table}

\textbf{MGnify}: Genus-level abundance tables were extracted from BIOM files via the MGnify REST API using a parallel collection pipeline (5 workers). \textbf{EMP}: Pre-processed closed-reference OTU table (Greengenes 13\_8) filtered to 4,628 soil samples by EMPO Level 3 classification. \textbf{Naylor et al.}: Raw 16S rRNA sequences (PRJNA369551) were processed using QIIME2 DADA2 pipeline in WSL Ubuntu, with taxonomy assigned via vsearch against SILVA 138 (99\% identity). \textbf{BonaRes}: 20-year long-term field trial data with bacterial community composition at genus level.

\subsection{Preprocessing Pipeline}

A six-step preprocessing pipeline was applied: (1) taxonomy unification to GTDB r220, (2) Total Sum Scaling normalization, (3) sparsity filtering (genera in $<$1\% of samples removed), (4) metadata standardization using ENVO ontology, (5) batch effect annotation, and (6) MGM-compatible tokenization (abundance-ranked sequences, length 512).

\subsection{Model Architecture and Training}

Gaia uses the GPT-2 architecture (8 layers, 8 attention heads, 256 embedding dimension, $\sim$9.7M parameters). Starting from MGM pre-trained weights (trained on 260,000 samples across all biomes), we expanded the vocabulary from 9,669 to 12,916 genera and applied a two-step continual pre-training strategy:

\begin{enumerate}
\item \textbf{Embedding-only training} (15 epochs, lr=$10^{-3}$): Only the token embedding layer is updated while all transformer weights are frozen. This allows 2,970 new soil-specific genera to learn meaningful representations without disrupting pre-trained knowledge.
\item \textbf{Full fine-tuning} (10 epochs, lr=$5 \times 10^{-5}$): All parameters are unfrozen for low-learning-rate adaptation to soil-specific patterns.
\end{enumerate}

Training was performed on a single NVIDIA RTX 5060 (8GB) with FP16 mixed precision, completing in approximately 25 minutes.

\subsection{Benchmark Tasks}

Ten tasks were defined spanning classification, regression, temporal prediction, and functional analysis:

\begin{enumerate}
\item \textbf{Biome Classification}: Forest vs grassland (MGnify)
\item \textbf{Drought Detection}: Drought vs control (Naylor et al., real experimental data)
\item \textbf{Abundance Reconstruction}: Next-token prediction accuracy
\item \textbf{Tillage Classification}: Cultivator vs plough (BonaRes)
\item \textbf{Fertilization Classification}: Extensive vs intensive (BonaRes)
\item \textbf{pH Prediction}: Soil pH regression (BonaRes)
\item \textbf{Current Yield Prediction}: Same-year yield regression (BonaRes)
\item \textbf{Future Yield Prediction}: This year's microbiome $\rightarrow$ next year's yield (BonaRes)
\item \textbf{Carbon Prediction}: Total soil carbon estimation (BonaRes)
\item \textbf{Functional Profile Analysis}: Drought vs control functional differences (Naylor)
\end{enumerate}

For tasks 1--9, a task-specific head (256$\rightarrow$128$\rightarrow$n) was trained on mean-pooled hidden states with backbone frozen. Random Forest (200 trees) served as baseline.

\section{Results}

\subsection{Overall Performance}

\begin{table}[H]
\centering
\caption{Benchmark results: Gaia vs Random Forest across ten tasks}
\label{tab:results}
\begin{tabular}{llccc}
\toprule
Task & Metric & Gaia & RF & Winner \\
\midrule
Biome Classification & Accuracy & 98.6\% & 98.7\% & Tie \\
Drought Detection & Accuracy & 91.9\% & 92.0\% & Tie \\
Abundance Reconstruction & Top-10 Acc & 31.3\% & N/A & Gaia \\
Tillage Classification & Accuracy & 93.1\% & 100.0\% & RF \\
Fertilization Classification & Accuracy & 70.7\% & 93.1\% & RF \\
pH Prediction & $R^2$ & 0.947 & 0.977 & RF \\
Current Yield & $R^2$ & 0.873 & 0.926 & RF \\
\textbf{Future Yield} & $R^2$ & \textbf{0.541} & 0.504 & \textbf{Gaia} \\
Carbon Prediction & $R^2$ & 0.829 & 0.948 & RF \\
Functional Analysis & -- & \multicolumn{2}{c}{See Section 3.5} & -- \\
\bottomrule
\end{tabular}
\end{table}

\subsection{Data Scaling Effects}

Performance scales consistently with training data quantity (Figure~\ref{fig:scaling}). Biome classification improved from 66.7\% (3,624 samples) to 98.6\% (8,329 samples). pH prediction improved from $R^2$=0.878 to 0.947. This monotonic improvement suggests our results underestimate the approach's potential with larger datasets.

\begin{figure}[H]
\centering
\includegraphics[width=\textwidth]{fig1_data_scaling.png}
\caption{Data scaling effects on three key benchmarks. Dashed lines indicate Random Forest performance. Gaia performance improves consistently with additional training data, approaching or matching RF.}
\label{fig:scaling}
\end{figure}

\subsection{Benchmark Comparison}

\begin{figure}[H]
\centering
\includegraphics[width=\textwidth]{fig2_benchmark_comparison.png}
\caption{Gaia vs Random Forest across benchmark tasks. Gold border highlights future yield prediction where Gaia outperforms RF.}
\label{fig:benchmark}
\end{figure}

\subsection{Temporal Prediction}

The most notable result is temporal prediction: using this year's microbiome to predict next year's yield (Figure~\ref{fig:future}). Gaia ($R^2$=0.541) outperforms Random Forest ($R^2$=0.504) by 7.3\% relative improvement. This suggests that the foundation model's understanding of microbial community dynamics provides advantages for predicting future soil states---a capability with direct agricultural applications.

\begin{figure}[H]
\centering
\includegraphics[width=\textwidth]{fig4_future_yield.png}
\caption{Future yield prediction: this year's microbiome predicting next year's yield. Gaia (left) achieves $R^2$=0.541 vs Random Forest (right) $R^2$=0.504.}
\label{fig:future}
\end{figure}

\subsection{Soil Carbon Prediction}

Gaia predicts total soil carbon from microbiome composition with $R^2$=0.829. While Random Forest achieves higher accuracy ($R^2$=0.948), this result demonstrates that microbial community structure alone contains substantial information about soil carbon storage---relevant for carbon credit verification and climate change monitoring.

\subsection{Functional Profile Analysis}

Using literature-based functional annotation of microbial genera, we compared functional profiles between drought and control samples from the Naylor dataset (Table~\ref{tab:function}).

\begin{table}[H]
\centering
\caption{Functional group changes under drought stress (Naylor et al., n=623)}
\label{tab:function}
\begin{tabular}{lrrr}
\toprule
Functional Group & Control & Drought & Change \\
\midrule
Carbon Decomposition & 2,146 & 4,262 & +98.6\% \\
Pathogen Suppression & 2,737 & 4,581 & +67.4\% \\
Denitrification & 788 & 578 & -26.6\% \\
Organic Matter Decomposition & 3,842 & 3,319 & -13.6\% \\
Nitrogen Fixation & 900 & 779 & -13.4\% \\
Plant Growth Promotion & 3,394 & 3,128 & -7.8\% \\
\bottomrule
\end{tabular}
\end{table}

Drought conditions substantially increase carbon decomposition (+98.6\%) and pathogen suppression (+67.4\%) activity while decreasing nitrogen fixation (-13.4\%) and denitrification (-26.6\%). These findings are consistent with published drought-microbiome studies and suggest that microbiome-based monitoring could provide early warning of soil functional degradation.

\subsection{Soil-Specialization vs General Model}

On abundance reconstruction, Gaia achieves 31.3\% Top-10 accuracy compared to 22.6\% for the original MGM---a 38\% relative improvement, validating the benefit of soil-specific continual pre-training. This task cannot be performed by Random Forest, representing a unique capability of the foundation model approach.

\section{Discussion}

\subsection{Foundation Models vs Traditional ML}

Gaia demonstrates that soil-specialized foundation models are competitive with Random Forest across diverse prediction tasks. The key advantage emerges in temporal prediction ($R^2$=0.541 vs 0.504), where understanding microbial community relationships is essential. Additionally, foundation models offer unique capabilities---abundance reconstruction, synthetic community generation, and functional interpretation---that traditional ML cannot provide.

\subsection{Carbon and Climate Applications}

The ability to predict soil carbon ($R^2$=0.829) from microbiome data has implications for carbon credit markets. Current soil carbon measurement requires physical sampling and chemical analysis (\$50--100 per sample). Microbiome-based estimation could provide a complementary, scalable verification tool. The functional analysis showing 98.6\% increase in carbon decomposition under drought further suggests that microbiome monitoring could predict carbon storage vulnerability.

\subsection{Agricultural Applications}

Gaia's temporal yield prediction capability ($R^2$=0.541) enables farmers to estimate next year's yield from this year's soil microbiome analysis. Combined with pH prediction ($R^2$=0.947), carbon estimation ($R^2$=0.829), and drought detection (91.9\%), Gaia offers comprehensive soil health assessment from a single microbiome sequencing.

\subsection{Limitations}

\begin{enumerate}
\item Training data (8,329 samples) is small compared to MGM (260K), though performance scales consistently with data quantity.
\item Temporal evaluation is limited to 283 time-series pairs from one experimental site.
\item Functional analysis uses literature-based annotation rather than direct functional gene measurement (e.g., shotgun metagenomics).
\item Several tasks (tillage, fertilization) favor RF, likely due to 20-year controlled experimental design creating highly separable features.
\end{enumerate}

\subsection{Future Directions}

\begin{enumerate}
\item Scaling to 50,000+ samples through additional public databases and collaborations
\item Integration of multi-modal inputs (microbiome + soil chemistry + climate data)
\item Direct functional prediction using PICRUSt2 or shotgun metagenomics
\item Prospective validation on independent experimental sites
\item Development of a web-based diagnostic tool for agricultural applications
\end{enumerate}

\section{Conclusion}

Gaia validates the concept of soil-specialized microbiome foundation models. By matching Random Forest on classification tasks, exceeding it on temporal prediction, and providing unique capabilities in functional analysis and community generation, Gaia demonstrates that learned microbial ``language'' representations capture ecologically meaningful patterns. The consistent improvement with data scaling, combined with the ability to predict soil carbon, pH, yield, and functional profiles from a single microbiome analysis, suggests a clear path toward practical tools for precision agriculture, soil health monitoring, and carbon credit verification.

\section*{Data and Code Availability}

All code, data collection pipelines, trained model weights, and benchmark scripts are available at \url{https://github.com/Kimchikilla/Gaia} under Apache 2.0 license.

\bibliographystyle{plainnat}
\begin{thebibliography}{9}

\bibitem[MGM(2024)]{mgm2024}
HUST-NingKang-Lab.
\newblock MGM: A Large-Scale Pretrained Foundation Model for Microbiome Analyses.
\newblock \emph{Advanced Science}, 2024.

\bibitem[Thompson et al.(2017)]{emp2017}
Thompson, L.~R. et al.
\newblock A communal catalogue reveals Earth's multiscale microbial diversity.
\newblock \emph{Nature}, 551:457--463, 2017.

\bibitem[Naylor et al.(2017)]{naylor2017}
Naylor, D. et al.
\newblock Drought and host selection influence bacterial community dynamics in the grass root microbiome.
\newblock \emph{The ISME Journal}, 11:2691--2704, 2017.

\bibitem[Raab et al.(2025)]{bonares2025}
Raab, M. et al.
\newblock Two decades long-term field trial data on fertilization, tillage, and crop rotation focusing on soil microbes.
\newblock \emph{Scientific Data}, 2025.

\bibitem[Bolyen et al.(2019)]{qiime2019}
Bolyen, E. et al.
\newblock Reproducible, interactive, scalable and extensible microbiome data science using QIIME 2.
\newblock \emph{Nature Biotechnology}, 37:852--857, 2019.

\end{thebibliography}

\end{document}
